% Обзор литературы. Уровни заголовков задаются \RWsection/\RWsubsection:
% в literature_review.tex — section/subsection, в main.tex — subsection/subsubsection.
Обзор охватывает пять направлений: (1)~архитектуру Transformer и её
теоретические свойства; (2)~механистическую интерпретируемость, которая
объясняет, \emph{что} вычисляют отдельные слои; (3)~недавнюю дискуссию о
том, являются ли нейросети <<уже чистыми>>; (4)~методы эффективной
аппроксимации самовнимания; (5)~сжатие и разреженность слоёв MLP. В конце
сформулированы выводы для нашей работы.

% =====================================================================
\RWsection{Архитектура Transformer и её свойства}

Архитектура Transformer~\cite{vaswani2017attention} отказалась от
рекуррентности и свёрток в пользу многоголового масштабированного
скалярного внимания $\mathrm{softmax}(QK^\top/\sqrt{d})V$, дополненного
позиционным кодированием, остаточными связями и поэлементными
слоями прямого распространения (MLP). Именно квадратичная по длине
последовательности стоимость внимания мотивирует методы из
раздела~\ref{sec:efficient}.

Два теоретических результата важны для понимания того, как устроены
слои. Dong и др.~\cite{dong2021attention} раскладывают выход сети из
одних слоёв внимания на сумму путей и доказывают, что без остаточных
связей и MLP такая сеть дважды экспоненциально по глубине сходится к
матрице ранга~1. Остаточные связи этому препятствуют, а MLP замедляет
вырождение. Отсюда следует, что остаточный поток "--- ключевой элемент
любой послойной декомпозиции. Cordonnier и
др.~\cite{cordonnier2020relationship} показывают, что многоголовое
самовнимание с относительным позиционным кодированием способно точно
выразить любой свёрточный слой (ядро $\sqrt{N_h}\times\sqrt{N_h}$ при
$N_h$ головах). Эксперименты на изображениях подтверждают, что
обученные головы действительно выучивают сеточные, свёрточные паттерны
вокруг текущей позиции.

\emph{Значение для нашей работы:} позиционные головы внимания можно
описывать фиксированными ядрами~\cite{cordonnier2020relationship}, а
заменять следует блок целиком, сохраняя остаточный
поток~\cite{dong2021attention}.

% =====================================================================
\RWsection{Механистическая интерпретируемость}

Программная работа <<Zoom In>>~\cite{olah2020zoom} формулирует три
гипотезы: признаки (features) "--- базовые единицы нейросетей;
признаки, связанные весами, образуют схемы (circuits); схожие признаки
и схемы возникают в разных моделях (универсальность). На примере
InceptionV1 разобраны детекторы кривых и другие схемы. Там же введено
понятие полисемантических нейронов и приведён пример, который позже
трактовался как суперпозиция.

Elhage и др.~\cite{elhage2021framework} переносят этот подход на
трансформеры. Остаточный поток рассматривается как линейный канал
связи, головы внимания "--- как независимые аддитивные компоненты с
QK-схемой (куда смотреть) и OV-схемой (что переносить). Для моделей
только из внимания получены точные описания: модель без слоёв
реализует статистику биграмм, однослойная "--- биграммы и
<<skip-триграммы>>, а двухслойная за счёт композиции голов строит
индукционные головы. Olsson и др.~\cite{olsson2022induction}
показывают, что индукционные головы (паттерн
$[A][B]\dots[A]\to[B]$) появляются в ходе обучения скачкообразно, в
<<фазовом переходе>>, совпадающем с резким улучшением обучения в
контексте. Для малых моделей связь доказана причинно, для больших
приведены корреляционные свидетельства.

Geva и др.~\cite{geva2021ffn} интерпретируют слой MLP как память
<<ключ--значение>>: ключи (строки первой матрицы) срабатывают на
текстовые шаблоны, а значения задают распределения над словарём
следующего токена. Принципиально для нас, что ключи \emph{нижних}
слоёв реагируют на поверхностные шаблоны (например, конкретные
$n$-граммы), а верхних "--- на семантические.

Bricken и др.~\cite{bricken2023monosemanticity} обучают разреженный
автокодировщик (словарное обучение) на активациях MLP однослойного
трансформера. Получаются тысячи признаков, значительно более
моносемантичных, чем исходные нейроны (например, признаки ДНК,
арабского письма, base64), а с ростом словаря признаки
<<расщепляются>> на более узкие. Rai и др.~\cite{rai2024review}
дают практический обзор области: таксономию задач (изучение
признаков, схем, универсальности), методов (probing, патчинг
активаций, словарное обучение и др.) и открытых проблем.

\emph{Значение для нашей работы:} есть точные описания вычислений
для моделей глубины 0--1~\cite{elhage2021framework}, свидетельства
<<поверхностности>> ранних MLP~\cite{geva2021ffn} и естественная
граница сложности "--- композиция слоёв, порождающая индукционные
головы~\cite{olsson2022induction}.

% =====================================================================
\RWsection{Дискуссия <<Neural networks are already clean>>}

В эссе Cammarata~\cite{cammarata_clean} (Goodfire, 2026; автор
называет его <<рабочей гипотезой>>) утверждается, что каноническое
свидетельство полисемантичности и суперпозиции было неверно
истолковано. Неинтерпретируемость поздних слоёв InceptionV1 объясняется
не суперпозицией, а архитектурными проблемами и слабыми
инструментами. После обычных улучшений обучения 99--100\% нейронов
последнего слоя становятся интерпретируемыми. Полисемантичность в
основном объясняется мотивом <<рекурсивного обобщения>>, то есть
объединением \emph{связанных} понятий. По предварительной оценке для
Qwen3-32B около 90\% <<полюсов>> нейронов имеют одну основную функцию,
около 7\% "--- две связанные, и лишь 2--4\% пока не объяснены. Вывод
автора: сеть можно изучать напрямую, без промежуточной декомпозиции
вроде разреженных автокодировщиков~\cite{bricken2023monosemanticity}.

Gao~\cite{gao_notclean} принимает главное наблюдение (полисемантичность
переоценивали), но возражает, что совпадение активирующих примеров с
шаблоном недостаточно для объяснения нейрона. Во-первых, из того, что
все активирующие примеры подходят под шаблон, не следует, что каждый
подходящий пример активирует нейрон. Во-вторых, важна величина
активаций: по его данным, обнуление слабых активаций заметно ухудшает
модель. Gao предлагает критерий: объяснение верно, если подстановка
\emph{смоделированных} по нему активаций вместо настоящих не
увеличивает KL-дивергенцию выходного распределения модели на большом
корпусе. Кроме того, понимание отдельных нейронов не даёт понимания
вычислений \emph{между} ними: найденные схемы остаются большими и
избыточными.

Обзорная заметка IBM Think~\cite{ibm_clean} пересказывает эссе и
собирает мнения исследователей: разметка нейронов не равна пониманию
модели; пики активаций покрывают малую долю срабатываний; словарное
обучение остаётся полезным для слабых активаций; нужны
интервенционные эксперименты. Работа Lissack~\cite{lissack_clean},
судя по названию (<<Already Clean, Not Yet Meaningful>>), разводит
<<чистоту>> внутренних механизмов и их содержательную осмысленность.

\emph{Значение для нашей работы:} задача замены первых слоёв явным
алгоритмом "--- это конструктивная проверка тезиса о <<чистоте>> на
уровне целых слоёв. Критерий Gao (downstream KL при подстановке)
совпадает с нашей основной метрикой качества аппроксимации, а его
замечание о слабых активациях указывает на главный риск табличных
аппроксимаций.

% =====================================================================
\RWsection{Эффективные аппроксимации внимания}\label{sec:efficient}

\RWsubsection{Ядерные и линейные методы}
Katharopoulos и др.~\cite{katharopoulos2020rnn} заменяют
$\exp(q^\top k)$ на скалярное произведение признаковых отображений
$\phi(q)^\top\phi(k)$ с $\phi(x)=\mathrm{elu}(x)+1$. Ассоциативность
матричного произведения даёт сложность $O(n)$, а авторегрессионная
модель становится рекуррентной сетью с постоянной памятью на шаг.
Performer~\cite{choromanski2021performer} (механизм FAVOR+)
строит \emph{несмещённую} оценку softmax-ядра положительными
ортогональными случайными признаками и доказывает оценки её
точности. RFA~\cite{peng2021rfa} аппроксимирует softmax-ядро
случайными фурье-признаками и добавляет механизм гейтинга для
смещения внимания к недавним токенам. SOFT~\cite{lu2021soft} отказывается
от softmax в пользу гауссова ядра, что позволяет построить
детерминированное низкоранговое разложение с опорными точками.
Псевдообратная матрица при этом вычисляется итерациями
Ньютона--Рафсона.

\RWsubsection{Низкоранговые методы}
Linformer~\cite{wang2020linformer} опирается на наблюдение, что
матрица внимания имеет низкий эффективный ранг, и проецирует $K$ и
$V$ по оси последовательности обучаемыми матрицами в размерность
$k\ll n$. Nyströmformer~\cite{xiong2021nystrom} аппроксимирует
softmax-матрицу методом Нистрёма с опорными точками, выбранными как
средние по сегментам последовательности, то есть
детерминированно. A3~\cite{wong2025a3} "--- метод сжатия \emph{после
обучения}: слой разбивается на функциональные компоненты QK, OV и MLP,
и для каждого аналитически находится уменьшение скрытой размерности,
минимизирующее функциональную ошибку компонента. В отличие от
разложения весовой матрицы на два сомножителя, это не добавляет
накладных расходов при выводе. На LLaMA~3.1-70B метод заметно
превосходит предыдущие низкоранговые подходы по perplexity при равном
бюджете.

\RWsubsection{Разреженные шаблоны внимания}
Sparse Transformer~\cite{child2019sparse} использует фиксированные
факторизованные (strided и fixed) шаблоны со сложностью
$O(n\sqrt{n})$ и моделирует последовательности длиной в десятки тысяч
элементов. Reformer~\cite{kitaev2020reformer} группирует близкие
запросы и ключи локально-чувствительным хешированием (случайные
повороты), получая $O(n\log n)$, и экономит память обратимыми
остаточными слоями. BigBird~\cite{zaheer2020bigbird} комбинирует
локальное окно, глобальные токены и случайные рёбра и доказывает, что
такое разреженное внимание сохраняет универсальность аппроксимации и
полноту по Тьюрингу.

\RWsubsection{Замена внимания}
FNet~\cite{leethorp2022fnet} заменяет подслой самовнимания
непараметрическим двумерным дискретным преобразованием Фурье и
достигает 92--97\% точности BERT на GLUE при существенно более
быстром обучении. Synthesizer~\cite{tay2021synthesizer} показывает,
что веса внимания можно синтезировать без попарного взаимодействия
токенов: по каждому токену в отдельности (Dense) или вообще
независимо от входа (Random). Такие модели остаются
конкурентоспособными на ряде задач.

\RWsubsection{Сводка}

\begin{table}[h]
\centering
\small
\begin{tabularx}{\textwidth}{@{}l X c c@{}}
\toprule
Метод & Способ аппроксимации & Случайность & Нужно обучение \\
\midrule
Linear Attention~\cite{katharopoulos2020rnn} & ядро $\mathrm{elu}+1$ & нет & да \\
Performer~\cite{choromanski2021performer} & случайные ортогональные признаки & да & да \\
RFA~\cite{peng2021rfa} & случайные фурье-признаки & да & да \\
SOFT~\cite{lu2021soft} & гауссово ядро + низкий ранг & нет & да \\
Linformer~\cite{wang2020linformer} & обучаемые проекции по длине & нет & да \\
Nyströmformer~\cite{xiong2021nystrom} & метод Нистрёма, сегментные средние & нет & да \\
A3~\cite{wong2025a3} & аналитическое снижение ранга QK/OV/MLP & нет & \textbf{нет} \\
Sparse Transformer~\cite{child2019sparse} & фиксированные разреженные шаблоны & нет & да \\
Reformer~\cite{kitaev2020reformer} & LSH-хеширование & да & да \\
BigBird~\cite{zaheer2020bigbird} & окно + глобальные + случайные рёбра & да & да \\
FNet~\cite{leethorp2022fnet} & преобразование Фурье & нет & да \\
Synthesizer~\cite{tay2021synthesizer} & синтезированные веса внимания & нет & да \\
\bottomrule
\end{tabularx}
\caption{Методы эффективного внимания. <<Нужно обучение>>: метод
предполагает обучение модели с нуля или дообучение, а не
построение по весам готовой модели.}
\end{table}

% =====================================================================
\RWsection{Сжатие и разреженность слоёв MLP}\label{sec:mlp}

Слой MLP в современных LLM (в Qwen3 "--- SwiGLU,
$\mathrm{down}(\mathrm{silu}(W_g x)\odot W_u x)$) содержит большую часть
параметров блока. Для его ускорения после обучения есть два подхода.

\paragraph{Низкоранговые разложения с учётом активаций.}
Простое SVD весовой матрицы плохо работает для LLM из-за выбросов в
активациях. ASVD~\cite{yuan2023asvd} масштабирует столбцы весовой
матрицы по статистике активаций, так что выбросы <<поглощаются>> весами
перед разложением, и подбирает ранг послойно с учётом
чувствительности слоёв. Без обучения это сжимает сеть на 10--30\%.
SVD-LLM~\cite{wang2025svdllm} предлагает <<отбеливание>> данных,
согласованное с усечением: после него каждое отброшенное сингулярное
число напрямую соответствует потере на калибровочных данных. Затем
сжатые веса дообучаются последовательной низкоранговой
аппроксимацией. Первая часть "--- замкнутая форма, вторая требует
градиентного обучения. Эти методы применимы к любым линейным
проекциям, включая проекции внимания. Отбор нейронов MLP в
A3~\cite{wong2025a3} относится к той же группе: это детерминированный
критерий в замкнутой форме.

\paragraph{Разреженность активаций.}
MoEfication~\cite{zhang2022moefication} исходит из наблюдения, что на
каждом входе активна малая доля нейронов FFN. Метод разбивает нейроны
на группы-<<эксперты>> и строит маршрутизатор, выбирающий группы для
входа. В итоге используется 10--30\% параметров FFN при сохранении
более 95\% исходного качества. Deja Vu~\cite{liu2023dejavu} обобщает
идею до <<контекстной разреженности>>: для каждого входа существует
малое подмножество голов внимания и нейронов MLP, дающее почти тот же
выход. Это подмножество предсказывается лёгким обучаемым предиктором
и используется для ускорения вывода больших моделей более чем в два раза.
Обе работы опираются на интерпретацию MLP как памяти
<<ключ--значение>>~\cite{geva2021ffn}: на конкретном входе срабатывают
немногие ключи.

\emph{Значение для нашей работы:} для MLP первых слоёв сравниваем
детерминированные части этих методов: разложения ASVD и SVD-LLM без
дообучения, разбиение MoEfication с маршрутизацией без обучения,
<<оракульную>> контекстную разреженность как верхнюю границу для Deja
Vu. Их обучаемые компоненты (дообучение SVD-LLM, предиктор Deja Vu)
оцениваем только в отдельном треке с дистилляцией.

% =====================================================================
\RWsection{Выводы и место нашей работы}

\begin{enumerate}
  \item Методы эффективного внимания применяют \emph{одну и ту же}
    аппроксимацию во всех слоях и, за исключением A3~\cite{wong2025a3},
    требуют обучения. Зависимость качества аппроксимации от глубины
    слоя в них не исследуется.
  \item Значительная часть методов стохастична
    (\cite{choromanski2021performer,peng2021rfa,kitaev2020reformer,zaheer2020bigbird}).
    Детерминированные альтернативы
    (\cite{xiong2021nystrom,lu2021soft,child2019sparse,leethorp2022fnet})
    существуют, но не сравнивались с ними именно как post-hoc
    замена слоёв обученной модели.
  \item Интерпретируемость даёт основания считать первые слои
    простыми: биграммы и skip-триграммы~\cite{elhage2021framework},
    поверхностные шаблоны в MLP~\cite{geva2021ffn}, свёрточные
    позиционные паттерны~\cite{cordonnier2020relationship}. Однако
    эти наблюдения не превращены в работающий алгоритм, который
    заменял бы слои целиком.
  \item Методы сжатия MLP~\cite{yuan2023asvd,wang2025svdllm,zhang2022moefication,liu2023dejavu}
    оценивают качество на всей модели при равномерном сжатии слоёв. Вопрос,
    насколько полно MLP \emph{первых} слоёв заменяется таблицами и
    линейными картами в замкнутой форме, остаётся открытым.
  \item Разреженность активаций без обучения
    (CATS~\cite{lee2024cats}, TEAL~\cite{liu2024teal}) отбирает нейроны
    порогом по величине активации и оценивается на всей модели; её
    поведение в первых слоях отдельно не изучалось. Удаление блоков с
    линейной заменой на границе (ShortGPT-подобный выбор блока,
    ReplaceMe~\cite{shopkhoev2025replaceme}, Ghosted
    Layers~\cite{yun2026ghosted}) использует регрессию на остаточном
    потоке "--- ту же конструкцию, что и наши калибровки, и потому
    подвержено той же ловушке.
  \item Массивные активации~\cite{sun2024massive} и
    <<супервеса>>~\cite{yu2024superweights} описаны как свойство
    обученных моделей и как источник трудностей квантования; их влияние
    на статистики пост-хок калибровки структурных замен и на отбор
    нейронов ранних слоёв до сих пор не рассматривалось.
  \item Дискуссия о <<чистоте>>
    нейросетей~\cite{cammarata_clean,gao_notclean} даёт и
    мотивацию, и критерий проверки: явная замена слоёв, оцениваемая по
    downstream KL, "--- строгая форма утверждения <<мы понимаем, что
    делают эти слои>>.
\end{enumerate}

Наша работа закрывает этот пробел: мы строим детерминированные
алгоритмы, заменяющие первые $k$ слоёв обученной LLM без дообучения,
и определяем глубину $k^{*}$, до которой такая замена возможна.
